\subsection{Modeling the activity landscape for cell-type-selective CRE design}

The public MPRA table contained 798,064 sequence records. After fixed quality
filtering, sequence validation, length filtering, and exact-sequence deduplication,
the full-data loader retained 749,054 records: 627,661 training, 58,811 validation,
and 62,582 test records. Chromosomes 19, 21, and X formed validation; chromosomes 7
and 13 formed test; remaining eligible chromosomes formed training. Parent selection
used identifiers and sequence hashes, never activity labels.

Forward/reverse-complement averaged Malinois predictions reproduced the published
test performance (mean Pearson $r=0.884$ and mean Spearman $\rho=0.825$ across the
three cell types). On the same 62,582 test records, the three-member residual-dilated
reviewer ensemble achieved mean Pearson $r=0.849$, whereas the three-member
multi-kernel evaluator achieved $r=0.793$ (Fig.~\ref{fig:validation}A). Their different
convolutional structures allowed one ensemble to review candidates during search and
the other to evaluate the frozen candidate set.

Predictive uncertainty was formed from learned heteroscedastic variance and
between-seed variation, with multiplicative scaling estimated exclusively on
validation data. Across cell types, 95\% test-interval coverage was 0.930--0.940 for
the reviewer and 0.923--0.935 for the cross-model evaluator
(Fig.~\ref{fig:validation}B). Reviewer uncertainty correlated with absolute error at
Spearman $\rho=0.399$--0.448; cross-model-evaluator uncertainty gave
$\rho=0.400$--0.467 (Fig.~\ref{fig:validation}C). Rejection curves showed decreasing
mean absolute error as high-uncertainty records were removed (Fig.~\ref{fig:validation}D),
supporting uncertainty-guided candidate triage.

\begin{figure}[t]
  \centering
  \includegraphics[width=\textwidth]{fig2_validation.pdf}
  \caption{\textbf{Predictive performance and uncertainty diagnostics.}
  \textbf{A}, test-set Pearson correlation by cell type for the primary, reviewer,
  and cross-model evaluator models. \textbf{B}, observed test coverage at 80\%, 90\%,
  and 95\% nominal levels after validation-only variance scaling; the dashed line
  marks ideal calibration. \textbf{C}, Spearman association between predicted
  uncertainty and absolute error. \textbf{D}, mean test error after progressively
  rejecting the most uncertain records. All neural ensembles contain three
  independently seeded members.}
  \label{fig:validation}
\end{figure}

\subsection{SafeEdit-CRE improves cross-model cell-type selectivity across 600 held-out parents}

We evaluated four editing methods on 600 held-out 200-nt natural CRE parents selected
by SHA-256 ordering after removal of all historical parent identifiers and sequence
hashes. Methods were applied for each of three target cells (K562, HepG2, SK-N-SH) and
four substitution budgets (1, 5, 10, 20 nt), producing 28,800 method-level comparisons.
Random editing used three independent replicates whose metrics were averaged without
best-of-three selection. After candidate generation, sequences were passed to the
held-out cross-model evaluator that had no role in the search.

Across all 7,200 parent--target--budget conditions, SafeEdit-CRE achieved mean
cross-model specificity-margin gain of 0.877, compared with 0.822 for iterative greedy
optimization (paired difference 0.055; parent-cluster bootstrap 95\% CI,
0.040--0.071) and 0.001 for random matched editing (difference 0.876; 95\% CI,
0.856--0.896) (Table~\ref{tab:g8}; Fig.~\ref{fig:primary}A). The fraction of
conditions with positive cross-model margin gain was 95.4\% for SafeEdit-CRE and
97.1\% for greedy, indicating that the advantage over greedy was driven by the
magnitude and distribution of gains rather than by a larger fraction of positive
cases. SafeEdit-CRE's advantage therefore came from larger gains across the
condition set, not from increasing the positive-case fraction.

The compute-matched primary beam search---which uses the same beam width, expansion
schedule, and search budget as SafeEdit-CRE but relies solely on the primary Malinois
predictor for candidate ranking---achieved a mean cross-model margin gain of 0.881.
SafeEdit-CRE matched this performance closely (paired difference $-0.0036$; 95\% CI,
$-0.0114$--0.0043; Fig.~\ref{fig:primary}B) while reducing absolute target gain by
0.205 (95\% CI, $-0.229$ to $-0.181$) and lowering cross-model ensemble uncertainty
by 0.0209 (95\% CI, $-0.0251$ to $-0.0166$) (Fig.~\ref{fig:primary}C--D). Relative to
greedy optimization, SafeEdit-CRE improved specificity margin by 0.055 (95\% CI,
0.040--0.071), reduced target gain by 0.141 (95\% CI, $-0.177$ to $-0.106$), and
reduced uncertainty by 0.0153 (95\% CI, $-0.0214$ to $-0.0091$).

\begin{table}[t]
\centering
\caption{Cross-model evaluation outcomes across 600 held-out parents.}
\label{tab:g8}
\resizebox{\textwidth}{!}{%
\begin{tabular}{lrrrr}
\toprule
Method & Mean margin gain & Mean target gain & Positive margin (\%) & Mean uncertainty \\
\midrule
Random matched & 0.001 & 0.012 & 48.6 & 0.185 \\
Greedy Malinois & 0.822 & 1.803 & 97.1 & 0.329 \\
Primary beam (compute-matched) & 0.881 & 1.866 & 94.8 & 0.335 \\
SafeEdit-CRE & 0.877 & 1.661 & 95.4 & 0.314 \\
\bottomrule
\end{tabular}
}
\end{table}

The compute-matched beam shows that search breadth alone can recover the same
predicted specificity. SafeEdit-CRE changes which member of that broad candidate pool
is selected: its variants have lower ensemble uncertainty and smaller target-score
excursions while retaining the specificity margin.

\begin{figure}[t]
  \centering
  \includegraphics[width=\textwidth]{fig3_primary_endpoint.pdf}
  \caption{\textbf{Cross-model evaluation on 600 held-out natural CRE parents.}
  \textbf{A}, mean held-out cross-model specificity-margin gain by editing
  method. \textbf{B}, paired difference in margin gain for SafeEdit-CRE versus
  greedy and the compute-matched primary beam. \textbf{C}, corresponding paired
  difference in absolute target gain. \textbf{D}, corresponding paired difference in
  cross-model ensemble uncertainty. Points show means; intervals are 100,000-replicate
  parent-cluster bootstrap 95\% confidence intervals. Random replicates were
  averaged, never selected by best outcome.}
  \label{fig:primary}
\end{figure}

The sequence-domain prefilter returned a zero-edit outcome for 12 of 600 parents
(2.0\%) in at least one condition. Retaining these outcomes in all summaries
preserved the prespecified denominator and quantified the practical cost of sequence
safeguards.

\subsection{Design-criteria benchmark on 90 non-overlapping confirmatory parents}

To assess whether SafeEdit-CRE designs not only score well under cross-model
evaluation but also satisfy a broad set of sequence-quality and model-agreement
criteria, we applied a prespecified nine-criterion design benchmark to 90 parent
sequences that did not overlap with any development or primary-evaluation parents
(six nominal parents were excluded after identifier-level overlap with an exploratory
benchmark was discovered).
For each method--parent--target--budget combination we recorded: non-negative primary
target gain, positive primary margin gain, positive reviewer margin transfer, strand
disagreement within domain, reviewer uncertainty within domain, naturalness within
domain, GC change within domain, homopolymer safety, and exact edit budget.
Thresholds were derived from validation-set random edits and locked before
evaluation.

SafeEdit-CRE passed all nine criteria for 654 of 1,080 designs (60.6\%), compared
with 419 of 1,080 greedy designs (38.8\%) and 142 of 1,080 random designs (13.1\%)
(Table~\ref{tab:g4}; Fig.~\ref{fig:hardaudit}A). The paired SafeEdit--greedy
difference was 21.8 percentage points (parent-cluster bootstrap 95\% CI,
18.3--25.3 percentage points). The SafeEdit--random difference was 47.4 percentage
points (95\% CI, 43.0--51.8 percentage points). The advantage over greedy was small
at a one-base budget and increased at budgets 5--20 (Fig.~\ref{fig:hardaudit}B--D).

\begin{table}[t]
\centering
\caption{Nine-criterion design benchmark across 90 non-overlapping confirmatory parents.}
\label{tab:g4}
\begin{tabular}{lrrrr}
\toprule
Method & Designs & Passed & Pass rate (\%) & Parent-bootstrap 95\% CI (\%) \\
\midrule
Random matched & 1,080 & 142 & 13.1 & 10.6--15.8 \\
Greedy Malinois & 1,080 & 419 & 38.8 & 34.6--42.9 \\
SafeEdit-CRE & 1,080 & 654 & 60.6 & 55.6--65.3 \\
\bottomrule
\end{tabular}
\end{table}

\begin{figure}[t]
  \centering
  \includegraphics[width=\textwidth]{fig3_hard_audit.pdf}
  \caption{\textbf{Nine-criterion design benchmark.} \textbf{A}, pooled nine-criterion
  pass rates in the locked non-overlapping 90-parent confirmatory benchmark. Error bars are parent-cluster bootstrap
  95\% confidence intervals. \textbf{B--D}, pass rates stratified by edit budget
  for K562, HepG2, and SK-N-SH, respectively.}
  \label{fig:hardaudit}
\end{figure}

The 21.8 percentage-point difference measures the benefit of enforcing reviewer,
uncertainty, and sequence-domain criteria during candidate generation. Post hoc
screening of greedy designs recovered markedly fewer variants that met the same
prespecified criteria.

\subsection{Cell-type and edit-budget effects}

The benefit of cross-model selection varied with target cell type and edit budget
(Fig.~\ref{fig:budgettradeoff}). For K562, the SafeEdit--greedy cross-model margin
difference progressed from $-0.016$ at one edit to 0.242 at 20 edits; for HepG2,
it progressed from $-0.004$ to 0.294. In SK-N-SH, differences ranged from $-0.078$
to $-0.033$ across budgets, indicating cell-type-specific heterogeneity in the
activity landscape.

Against the compute-matched primary beam, SafeEdit-CRE differences remained within
tight equivalence bounds in all strata: K562 0.004--0.031; HepG2 $-0.030$--0.004;
SK-N-SH $-0.033$--0.010, confirming that the pooled margin equivalence holds across
target cells and edit budgets.

\begin{figure}[t]
  \centering
  \includegraphics[width=\textwidth]{fig5_budget_tradeoff.pdf}
  \caption{\textbf{Cell-type and edit-budget effects.} \textbf{A--B}, cross-model
  specificity-margin gain and uncertainty pooled over the three target cells across
  edit budgets for greedy, primary-beam, and SafeEdit-CRE searches. \textbf{C--D},
  target-by-budget differences in mean specificity-margin gain for SafeEdit-CRE
  relative to greedy or the compute-matched primary beam. Random editing is omitted
  from panels A--B because its pooled margin gain remains near zero. Lines and
  heatmap cells are descriptive means over parents. Edit budgets are reported in nt
  (nucleotide substitutions).}
  \label{fig:budgettradeoff}
\end{figure}

The pattern is consistent with consensus reranking becoming increasingly valuable as
the number of plausible edit paths grows. At one substitution, all three directed
strategies are frequently constrained to the same small neighborhood of the parent
sequence; at 5--20 substitutions, cross-model review filters out edits that look
favorable under the primary predictor but transfer poorly to the held-out
ensembles. The SK-N-SH result suggests that design thresholds should be calibrated
per cell type rather than transferred unchanged across cellular models.

\subsection{Complete re-search ablation identifies key components}

To identify which SafeEdit-CRE components drive the observed behavior, we reran the
complete search under seven configurations on 180 parents (15,120 rows). Each
configuration generated its own candidate sequences; filters were not removed
retrospectively from a shared candidate table. Cross-model inference was recomputed
for every candidate after each search.

Full SafeEdit-CRE yielded mean cross-model margin gain 0.890, target gain 1.705,
and uncertainty 0.312 (Fig.~\ref{fig:ablation}). The matched primary beam had margin
gain 0.890 (paired full-minus-primary difference 0.0002; 95\% CI, $-0.0145$--0.0151),
target gain 1.897 (difference $-0.192$; 95\% CI, $-0.244$ to $-0.143$), and
uncertainty 0.334 (difference $-0.0212$; 95\% CI, $-0.0283$ to $-0.0140$).

Removing naturalness reranking raised cross-model margin gain to 0.979 and target
gain to 2.017, together with an uncertainty increase to 0.343. Full SafeEdit-CRE
reduced uncertainty by 0.0305 (95\% CI, 0.0218--0.0394) relative to this
configuration while moderating margin by 0.089 (95\% CI, 0.075--0.102) and target
gain by 0.312 (95\% CI, 0.266--0.360). Removing the sequence-domain prefilter
eliminated all 48 infeasible rows and raised margin gain to 0.910 by relaxing the GC
and homopolymer limits. Strand and uncertainty penalties had smaller pooled effects,
whereas the reviewer-free configuration reproduced the primary beam exactly.

\begin{figure}[t]
  \centering
  \includegraphics[width=\textwidth]{fig6_ablation.pdf}
  \caption{\textbf{Complete re-search ablation on 180 parents.} \textbf{A}, change in
  cross-model specificity-margin gain relative to full SafeEdit-CRE. \textbf{B},
  corresponding change in absolute target gain. \textbf{C}, mean cross-model
  ensemble uncertainty for selected candidates. \textbf{D}, number of infeasible
  search rows. Each configuration was independently rerun. Naturalness-aware
  reranking is the primary driver of uncertainty reduction, while prefilter
  enforcement trades margin for synthesis feasibility.}
  \label{fig:ablation}
\end{figure}

Naturalness-aware reranking made the largest contribution to uncertainty reduction,
whereas the sequence-domain prefilter determined synthesis feasibility. Removing the
reviewer reproduced the primary-beam result exactly, showing that the reviewer changes
the returned sequence rather than merely annotating it after selection.

\subsection{Tier A candidate library and representative examples}

From the 3,240 designs in the 90-parent confirmatory benchmark, we derived a tiered candidate
library for experimental screening. Seventy-four Tier A sequences were generated by
SafeEdit-CRE, passed all nine criteria, met target and margin requirements, had no
synthesis-risk flag, and lay on or near the primary--reviewer Pareto frontier
(Supplementary Table~S9; Supplementary Fig.~S3). Tier B contains 1,141 alternatives
that passed core activity and cross-model checks but were Pareto-dominated or mildly
outside a sequence-quality criterion; the remaining 2,025 Tier C sequences are
retained for failure-mode analysis. The resulting 74-sequence Tier A panel spans all
three target cell types and four edit budgets, providing a compact first-pass
reporter-screening set, while Tier B retains sufficient breadth for a larger secondary
MPRA library.

To illustrate the nature of SafeEdit-CRE designs, Supplementary Table~S10 presents
three representative Tier A candidates---one each for K562, HepG2, and SK-N-SH---all
generated under a 5-edit budget and achieving Pareto-frontier status without synthesis
risk flags. Each example reports the parent and edited 200-nt sequence, edit positions,
primary and reviewer specificity-margin gains, primary target-cell activity gain, the
reduction in cross-model predictive uncertainty, GC content, naturalness, and
homopolymer status. All three show positive primary target gain and positive reviewer
margin transfer with decreased cross-model uncertainty. Local sequence changes near
edited positions are recorded as testable sequence-grammar hypotheses for experimental
follow-up.
